\subsection{Actual intrinsic finite regular-unit convex PL algebra}

\begin{lemma}[The same original section units restrict and concatenate]
\label{scheme:actual-native-unit-expression-algebra}
\leanok
\lean{EnvelopeProperty.Berkovich.RelativeSchemeValuedPoint.unitTropAffine_restrict,
EnvelopeProperty.Berkovich.RelativeSchemeValuedPoint.unitTropAffine_append,
EnvelopeProperty.Berkovich.RelativeSchemeValuedPoint.unitTropAffine_append_zero_right,
EnvelopeProperty.Berkovich.RelativeSchemeValuedPoint.unitTropAffine_zero_append}
\uses{scheme:native-regular-unit-local-psh,scheme:valued-point-topology}
For an ACTUAL inclusion $V'\subseteq V$ of scheme opens, the genuine original
restriction RingHom sends units of $\Gamma(X,V)$ to units of
$\Gamma(X,V')$. Every rational affine unit expression evaluated at an actual
native point in $V'$ equals the original expression of the original units.
For two finite families on the SAME true open, concatenating the genuine
units and rational coefficient vectors adds the two entire real expressions;
zero extension of one coefficient vector recovers that entire original
expression. No coefficient algebra action on a section ring is inferred.
\end{lemma}
\begin{proof}
Apply the true presheaf restriction RingHom to the original units and their
inverses. Actual native section evaluation commutes with this WHOLE restriction
on every original section element. Therefore each original logarithmic term,
its finite sum and its rational constant are unchanged. For concatenation
split the actual finite sum over the two original finite index sets. True
rational-cast arithmetic gives the original sum identity. Extending one
rational vector by zero makes every new term zero, recovering the entire
original expression. These are identities of original evaluations, never an
inference that selected values determine a seminorm.
\end{proof}

\begin{theorem}[Finite original convex PL algebra]
\label{scheme:actual-native-convex-pl-closure}
\leanok
\lean{EnvelopeProperty.Berkovich.RelativeSchemeValuedPoint.IsConvexPLOn.isOpen,
EnvelopeProperty.Berkovich.RelativeSchemeValuedPoint.IsConvexPLOn.mono,
EnvelopeProperty.Berkovich.RelativeSchemeValuedPoint.IsConvexPLOn.congr,
EnvelopeProperty.Berkovich.RelativeSchemeValuedPoint.isConvexPLOn_of_local,
EnvelopeProperty.Berkovich.RelativeSchemeValuedPoint.isConvexPLOn_const,
EnvelopeProperty.Berkovich.RelativeSchemeValuedPoint.IsConvexPLOn.max,
EnvelopeProperty.Berkovich.RelativeSchemeValuedPoint.IsConvexPLOn.add,
EnvelopeProperty.Berkovich.RelativeSchemeValuedPoint.IsConvexPLOn.smul,
EnvelopeProperty.Berkovich.RelativeSchemeValuedPoint.IsConvexPLOn.add_const,
EnvelopeProperty.Berkovich.RelativeSchemeValuedPoint.isConvexPLOn_finset_sum,
EnvelopeProperty.Berkovich.RelativeSchemeValuedPoint.isConvexPLOn_finset_sup'}
\uses{scheme:actual-native-unit-expression-algebra}
On ANY genuine prescribed-base native scheme space, intrinsic convex PL
functions form a local class, restrict to actual open subsets, and preserve
equality on their original domains. Finite nonempty maxima, finite sums,
nonnegative rational multiples and rational translates preserve this class.
Rational constants are convex PL on every actual open domain.
No field, finite type, Noetherianity, separation, completeness or norm bound
is needed. The original scheme structure and EVERY coefficient value are
retained.
\end{theorem}
\begin{proof}
Each definition supplies a true open neighborhood contained in its domain;
these cover the domain, proving it is open. Restriction intersects the original
neighborhood with the new genuine open. Equality changes only the original
function identity, and locality takes the actual presentation on the supplied
local neighborhood.

For two functions, intersect their true analytic neighborhoods and restrict
BOTH original section-unit families to the actual intersection of their two
scheme opens. The preceding whole restriction identities preserve both
original formulas there. Concatenate those actual unit families. For the
maximum, extend the two rational vector families by zero and take their true
finite union; its actual nonempty maximum is the maximum of the two original
maxima. For the sum, take the actual finite product of rational vector families
and concatenate each pair of vectors with the sum of its constants. The
maximum of all pairwise sums is exactly the sum of the two original finite
maxima, by genuine finite supremum distributivity for real addition. All
original finite real values embed into the actual extended-real operations.

For a nonnegative rational factor, multiply EVERY original rational vector
and constant by it. Multiplication by that same nonnegative real preserves
finite maxima, including the zero case. A rational constant uses the empty
unit family and a single original rational affine constant. Adding that
constant gives rational translation. Induction gives finite sums and finite
nonempty maxima; an empty sum is the genuine constant zero.
\end{proof}
